\section{컴파일 산출물과 적합성 조건}
\label{sec:conformance}
\subsection{실행 manifest}
컴파일러는 프로그램의 정규화된 구문, 형식화 그래프, 하드웨어 바인딩과 실행 계획을
하나의 manifest로 생성한다. manifest는 값의 표현과 의미를 함께 추적하는 기준 문서이다.
표~\ref{tab:manifest}는 그 필드 집합을 정의한다.

\begin{table}[tb]
\centering\small
\caption{실행 manifest의 필수 구성}
\begin{tabularx}{\linewidth}{@{}P{.23\linewidth}Y@{}}
\toprule 구성 & 저장 내용 \\\midrule
언어와 그래프 & 언어 버전, 정규화 AST, 그래프 식별자, 원시 연산자 버전 \\
하드웨어 & 프로필 $H$, 설정 $\theta$, 기능 요구, 접근 조건, 원문 식별자 \\
형식 & 채널 순서, 측정량, 단위, reference, 좌표계, 자료형, stride \\
시간 & 원래 시계, clock epoch, 변환 버전, 유효 구간, 오차 계약 \\
수치 & 계수, 표현 정밀도, 반올림ㆍ포화 정책, 참조 오차 기준 \\
학습 & 분할 manifest, fit 계보, 모델 digest, 모델 활성화 사건 \\
실행 & actor 소유권, 일정, 버퍼 용량, 큐 정책, 서비스ㆍ응답시간 계약 \\
기록 & 원본 범위, 품질 사건, 스키마 변경, 실행 환경과 결과 식별자 \\
\bottomrule
\end{tabularx}
\label{tab:manifest}
\end{table}

형식화 구체화는 단위ㆍ채널ㆍ시계 제약을 풀고 각 연결선의 타입을 결정한다.
정적 인덱스 제약은 정수 선형 산술, 유리수 비교와 유한 집합 관계로 구성한다.
행렬의 수치 조건과 장치 조회값은 해당 원시 연산자ㆍ프로필의 계약 식별자에 연결한다.
원시 연산자 선택은 수치 커널과 상태 표현을 결정한다. 실행 계획 생성은
고정률 부분그래프의 firing 순서, 사건 경계, 메모리 소유권과 버퍼 위치를 결정한다.
각 변환은 입력ㆍ출력 그래프의 대응표와 계약 증거를 manifest에 추가한다.

\subsection{관측 의미의 대응}
프로그램 $P$, 입력 trace $I$, 환경 $A$에서 코어 의미가 허용하는 실행 집합을
$\mathcal T(P,I,A)$로 둔다. 정수 디코딩, 사건 순서와 모델 버전은 정확한 관측 관계를
사용하고, 실수 결과는 식~\eqref{eq:numeric}의 수치 관계를 사용한다.
컴파일 결과 $C(P)$의 실행 trace $T$에 대해 적합성 관계를
\begin{equation}
 \operatorname{Obs}(T)
 \in\mathcal T_{\epsilon}(P,I,A)
 \label{eq:refinement}
\end{equation}
로 정의한다. $\mathcal T_\epsilon$은 명시한 수치 허용오차와 시간 구간을 적용한 관측 집합이다.
관측에는 값, 가용 시각, 품질, 모델 버전, schemaㆍclock epoch와 외부 사건을 포함한다.
버퍼 정책이 생성하는 gap과 상태 전이도 같은 trace의 원소로 기록한다.

검사기는 각 원시 연산자의 전이 불변식을 확인하고, 합성 규칙에 따라 출력 계약을 계산한다.
하드웨어 바인딩은 식~\eqref{eq:bind}를 평가한다. 일정 검사는 반복 벡터와 실행 접두의
토큰 수를 검사하고, 시간 검사는 선택한 arrivalㆍservice curve 또는 응답시간 계약을 계산한다.
각 결과는 적용한 입력 영역, 구성, 도구ㆍ계약 버전과 함께 저장한다.

\subsection{규격 적합성의 판정 대상}
\begin{longtable}{@{}P{.19\linewidth}P{.36\linewidth}P{.39\linewidth}@{}}
\caption{적합성 검사 대상과 판정 관계}\label{tab:conformance}\\
\toprule 대상 & 입력과 근거 & 판정 관계 \\\midrule
\endfirsthead
\multicolumn{3}{l}{\small 표 \thetable. 적합성 검사 대상과 판정 관계 (계속)}\\
\toprule 대상 & 입력과 근거 & 판정 관계 \\\midrule
\endhead
\bottomrule\endfoot
패킷 디코딩 & 프로토콜 판본, 패킷ㆍ기준값 쌍, 경계 길이 & 부호ㆍbyte orderㆍ채널ㆍ상태값의 정확한 대응 \\
물리량 변환 & 정수 범위, gain, offset, 기준 단위 & 기준값과 식으로 정의한 수치 관계 \\
시계 변환 & 동기 사건 쌍, 관측 경계, 유효 구간 & 시간 구간 포함과 오차 증거 종류 \\
스키마 전이 & 설정 전후 descriptor와 경계 사건 & 구간별 형식, 상태 변환, epoch 순서 \\
품질 전파 & gapㆍ보간ㆍ복구ㆍ포화 사건 trace & 연산자 정책에 따른 출력 마스크와 계보 \\
학습 계보 & 원천 ID, 그룹 함수, fitㆍtuneㆍeval 집합 & 식~\eqref{eq:split} 및 모델 가용 조건 \\
실행 일정 & 생산ㆍ소비율, 반복 벡터, 초기 토큰 & firing 가능성, 주기 평형, 최대 점유량 \\
시간ㆍ메모리 & 유입 상한, 서비스 곡선, 환경 조건 & backlogㆍ반응시간ㆍ메모리 예산 \\
기록ㆍ재생 & manifest, 원본ㆍ사건 순서, 활성 버전 & 식~\eqref{eq:refinement}의 관측 대응 \\
\end{longtable}

\subsection{설정 변경과 재연결의 전이}
설정 변경 사건은 새 구성 $\theta'$와 새 스키마 기술자를 생성한다.
실행기는 현재 epoch의 진행 경계를 닫고, 남은 처리 의존성을 해당 epoch에 귀속시킨다.
다음 epoch에는 새로운 형식과 호환하는 상태 변환 $R_Z$를 적용한다.
\begin{equation}
 R_Z:Z_{\theta}\longrightarrow Z_{\theta'},\qquad
 R_Z(\mathcal J_{\theta})\subseteq\mathcal J_{\theta'}.
\end{equation}
초기화 정책은 $R_Z(z)=z^0_{\theta'}$로 표현한다.
장치 재연결은 acquisition epoch를 갱신하고, 시계 재설정은 clock epoch와
새 시간 변환 레코드를 생성한다. 이 사건들이 배열 해석과 시간 연속성의 경계를 결정한다.

\subsection{기록과 재생의 관측 경계}
기록 어댑터는 원래 측정 시각, 원래 가용 시각과 재생 시각을 각각 저장한다.
실시간 실행 재생은 원래 가용 순서를 사용하여 $\cut_u$를 구성한다.
회고 분석은 새 시간 뷰와 변환 계보를 생성하여 기록 구간을 입력으로 사용한다.
XDF는 스트림 timestamp와 clock offset 정보를 기록하고\cite{xdf,lsl_time},
SNIRF는 광학 측정열ㆍ시간 배열과 chunked 기록을 표현한다\cite{snirf}.
Open Ephys와 SpikeGLX의 기록 어댑터는 해당 버전의 표본 번호ㆍ시각ㆍ스케일 필드를
복원한다\cite{open_ephys_binary,spikeglx_metadata}.
재생기는 이 기록 형식들을 공통 시간 레코드와 계보 그래프에 대응시킨다.
